\chapter{गति: एक-सी और बदलती हुई}\label{ch:g9:uniform-varied-motion}

दो साल पहले तुमने गति का निदान आँख से किया था: बराबर फ़ासलों पर बिंदु,
तो एक-सी गति; फैलते हुए बिंदु, तो तेज़ होती हुई। इस साल उन बिंदुओं को
मिलीमीटर के निशान मिलते हैं और निदान को अंक --- अंतराल दर अंतराल, \omterm{def:g7:motion-average-speed:formula}{चाल}
दर \omterm{def:g7:motion-average-speed:formula}{चाल}। और अध्याय के आख़िर में \omterm{def:g2:pushes-and-pulls:force}{बल} तथा गति के बारे में एक शांत-सा वाक्य
खड़ा है, जो एक दिन आधी भौतिकी अपनी पीठ पर ढोएगा।

\section{गति का अभिलेख}

\begin{definition}[समयबद्ध स्थिति-अभिलेख]\label{def:g9:uniform-varied-motion:record}
किसी गति का \emph{समयबद्ध स्थिति-अभिलेख}\index{समयबद्ध स्थिति-अभिलेख}
चलती वस्तु की जगहों को समय की बराबर टिक-टिक पर दर्ज करता है --- यानी
एक-एक करके देखे गए वीडियो के फ़्रेम, कौंधती स्ट्रोब से खींचा चित्र, या
प्रयोगशाला का बिंदु-कालमापी, जो वस्तु से खिंचती काग़ज़ की पट्टी पर हर
सेकंड पचास बिंदु ठोंकता है। किन्हीं दो पड़ोसी निशानों के बीच एक टिक भर
का सफ़र पड़ा होता है: यह अभिलेख गति को मापे जा सकने वाले टुकड़ों में
बदल देता है।
\end{definition}

\begin{method}[अभिलेख को अंकों के साथ पढ़ना]\label{met:g9:uniform-varied-motion:read}
किसी ऐसे अभिलेख के लिए जिसकी टिक की लंबाई $\tau$ हो (मान लो हर सेकंड
पचास के लिए \qty{0.02}{s}, या वीडियो से मिली आरामदेह \qty{0.1}{s}):
\begin{enumerate}
  \item पड़ोसी निशानों के बीच का हर फ़ासला मिलीमीटर वाले रूलर से
        मापो;
  \item हर फ़ासले को टिक से भाग दो: $v = d/\tau$ उस छोटे अंतराल पर
        (औसत) \omterm{def:g7:motion-average-speed:formula}{चाल} दे देता है --- और अगर फ़ासले \omterm{def:g2:measuring-length:units}{मीटर} में हों तो
        \unit{m/s} में;
  \item चालों को कंधे से कंधा मिलाकर रखो और कहानी पढ़ो: टिके हुए अंक,
        तो एक-सी गति; चढ़ते हुए, तो त्वरित; उतरते हुए, तो मंदित;
  \item किसी भी अंतराल की \omterm{def:g7:motion-average-speed:formula}{चाल} को उस पल ``की'' \omterm{def:g7:motion-average-speed:formula}{चाल} कहकर बताओ: इतनी छोटी
        टिक पर औसत और तात्क्षणिक आपस में हाथ मिला लेते हैं।
\end{enumerate}
\end{method}

\begin{omfigure}
\begin{tikzpicture}[scale=0.85]
  \foreach \x in {0.5,2.0,3.5,5.0,6.5,8.0,9.5} {
    \fill[omDef] (\x,2.6) circle (0.1);
  }
  \node[right, font=\small] at (10.0,2.6)
    {सारे फ़ासले \qty{3.0}{cm}: एक-सी गति};
  \foreach \x in {0.5,1.1,2.0,3.2,4.7,6.5,8.6} {
    \fill[omThm] (\x,1.5) circle (0.1);
  }
  \node[right, font=\small] at (10.0,1.5)
    {बढ़ते फ़ासले: त्वरित};
  \foreach \x in {0.5,2.6,4.3,5.6,6.5,7.0,7.2} {
    \fill[omMeth] (\x,0.4) circle (0.1);
  }
  \node[right, font=\small] at (10.0,0.4)
    {घटते फ़ासले: मंदित};
\end{tikzpicture}

{\small बिंदु-कालमापी से निकली तीन पट्टियाँ, हर टिक पर एक बिंदु। इस
साल ये फ़ासले मापे जाते हैं, और हर एक एक \omterm{def:g7:motion-average-speed:formula}{चाल} बन जाता है।}
\end{omfigure}

\begin{example}[एक पट्टी, हल करके]\label{ex:g9:uniform-varied-motion:worked}
किसी ठेले की पट्टी, \qty{0.1}{s} की टिकें, और \omterm{def:g2:measuring-length:units}{सेंटीमीटर} में फ़ासले:
$2.0$, $3.0$, $4.0$, $5.0$, $6.0$. फ़ासले दर फ़ासले चालें:
$0.02 \div 0.1 = \qty{0.2}{m/s}$, फिर $0.3$, $0.4$, $0.5$,
\qty{0.6}{m/s}. फ़ैसला: त्वरित --- और वह भी ख़ूबसूरती से नियमित ढंग
से: \omterm{def:g7:motion-average-speed:formula}{चाल} हर टिक पर ठीक \qty{0.1}{m/s} बढ़ती है। जिन गतियों की \omterm{def:g7:motion-average-speed:formula}{चाल} बराबर
समयों में बराबर क़दमों से चढ़ती है, उनका एक ख़ास नाम है और एक शानदार
भविष्य --- \emph{एक-समान} त्वरित, यानी टिके हुए कारणों के दस्तख़त।
\end{example}

\begin{example}[सबसे मशहूर त्वरित गति]\label{ex:g9:uniform-varied-motion:freefall}
किसी पत्थर को गिराकर उसका वीडियो बनाओ: फ़ासले टिक दर टिक खिंचते जाते
हैं --- मुक्त पतन त्वरित है। मापने पर उसकी \omterm{def:g7:motion-average-speed:formula}{चाल} गिरने के हर सेकंड
में क़रीब \qty{9.8}{m/s} बढ़ती है: एक सेकंड बाद \qty{9.8}{m/s}; और दो
के बाद लगभग \qty{20}{m/s}। इस बढ़त की दर का उस जगह की $g$ के बराबर
होना कोई इत्तेफ़ाक़ नहीं है --- यह इस साल की भौतिकी की सबसे गहरी तुक है,
और हाई स्कूल वाली जिल्द अपनी यांत्रिकी उसी पर खड़ी करती है। इसे दर्ज
करो, रेखांकित करो, और इंतज़ार करने दो।
\end{example}

\section{तात्क्षणिक चाल, ईमानदारी से}

\begin{definition}[तात्क्षणिक चाल]\label{def:g9:uniform-varied-motion:instant}
किसी पल पर किसी गति की \emph{तात्क्षणिक चाल}\index{तात्क्षणिक चाल}
उस पल के आस-पास के बहुत छोटे अंतराल पर की \omterm{def:g7:motion-average-speed:average}{औसत चाल} है --- इतने छोटे कि
उसके भीतर गति को अपनी रफ़्तार बदलने की जगह ही न मिले। यही गतिमापी का
अंक है और रडार का भी: दरअसल दोनों ठीक वैसे ही मापते हैं जैसे
\cref{met:g9:uniform-varied-motion:read}, बस पलक भर की टिकों पर।
(``बहुत छोटा'' का पूरी कठोरता के साथ मतलब क्या है, यही वह महान सवाल है
जिसका जवाब हाई स्कूल के गणित का आख़िरी साल देता है --- और उसी जवाब ने
आधुनिक विज्ञान गढ़ा।)
\end{definition}

\begin{example}[औसत और तात्क्षणिक, आमने-सामने]\label{ex:g9:uniform-varied-motion:sidebyside}
दो स्टेशनों के बीच मेट्रो की एक दौड़: \omterm{def:g7:motion-average-speed:average}{औसत चाल} --- कुल दूरी बटा कुल समय
--- \qty{35}{km/h}. \omterm{def:g9:uniform-varied-motion:instant}{तात्क्षणिक चाल}: दोनों प्लेटफ़ॉर्मों पर $0$, और
सुरंग के बीच \qty{70}{km/h}। एक-सी गति में, और सिर्फ़ वहीं, दोनों
विचार मिलकर एक ही टिके हुए अंक बन जाते हैं: एक-सी गति ठीक वही गति है
जो \omterm{def:g7:motion-average-speed:formula}{चाल} को एक अकेली, ईमानदार, अंतराल-मुक्त राशि बना देती है।
\end{example}

\section{वह शांत वाक्य}

\begin{proposition}[संतुलित बल गति को अनछुआ छोड़ देते हैं]\label{prop:g9:uniform-varied-motion:inertia}
जब किसी पिंड पर लगते \omterm{def:g2:pushes-and-pulls:force}{बल} एक-दूसरे को काट देते हैं --- या जब कोई \omterm{def:g2:pushes-and-pulls:force}{बल} लगता
ही नहीं --- तब उस पिंड की गति नहीं बदलती: ठहरा हुआ ठहरा ही रहता है, और
चलता हुआ सीधी रेखा में अटल \omterm{def:g7:motion-average-speed:formula}{चाल} से चलता रहता है। गति में बदलाव --- तेज़
होना, धीमा होना, मुड़ना --- सिर्फ़ तभी होता है जब कोई \emph{असंतुलित}
\omterm{def:g2:pushes-and-pulls:force}{बल} लग रहा हो। यह वाक्य, जिसे यहाँ कहकर ईमानदारी से आगे के लिए टाल दिया
जा रहा है, पूरी यांत्रिकी का पहला नियम है; हाई स्कूल वाली जिल्द इसी से
शुरू होती है और फिर इसे कभी नहीं छोड़ती।
\end{proposition}

\begin{example}[इस वाक्य से दुनिया पढ़ना]\label{ex:g9:uniform-varied-motion:reading}
चिकनी बर्फ़ पर सरकती पक सीधी और टिकी हुई फिसलती है --- \omterm{def:g2:pushes-and-pulls:force}{बल} \omterm{def:g2:balance-equilibrium:equilibrium}{संतुलित}
(\omterm{def:g4:weight-and-mass:weight}{भार} नीचे, बर्फ़ का सहारा ऊपर), गति अनछुई: चलते \emph{रहने} के लिए
किसी धक्के की ज़रूरत नहीं। ब्रेक लगाती कार धीमी होती है: एक असंतुलित
\omterm{def:g2:pushes-and-pulls:force}{बल}, यानी टायरों पर सड़क की पकड़, पीछे की ओर लगता हुआ। मुड़ता साइकिल
सवार: बग़ल की ओर असंतुलित \omterm{def:g2:pushes-and-pulls:force}{बल}, जो झुके हुए टायर देते हैं। और टिकी हुई
\qty{900}{km/h} पर उड़ता यात्री विमान: प्रणोद खिंचाव को काटता हुआ, और
उत्थापन \omterm{def:g4:weight-and-mass:weight}{भार} को --- यानी चार \omterm{def:g2:pushes-and-pulls:force}{बल}, ठीक बराबरी, और एक-सी गति। ठहराव और
सफ़र भौतिकी की नज़र में एक ही शांत हालत हैं।
\end{example}

\begin{remark}[यह वाक्य चौंकाता क्यों है]\label{rem:g9:uniform-varied-motion:surprise}
रोज़मर्रा की ज़िंदगी विरोध करती लगती है: पैडल मारना बंद करो और साइकिल
धीमी पड़ जाती है --- तो क्या गति को \omterm{def:g2:pushes-and-pulls:force}{बल} की \emph{ज़रूरत} नहीं? पर धीमी
पड़ती साइकिल बल-मुक्त है ही नहीं: घर्षण और \omterm{def:g2:air-around-us:air}{हवा} पीछे की ओर धकेलते हैं,
असंतुलित होकर, और ठीक इसीलिए रफ़्तार क्षय होती है। उन्हें हटा दो ---
बर्फ़ का मैदान, अंतरिक्ष का शून्य --- और गति हमेशा के लिए बहती रहती
है: तुम्हारे दादा-दादी के मिलने से भी पहले छोड़े गए अंतरिक्ष यान आज भी
ठंडे इंजनों के साथ बहते जा रहे हैं। घर्षण का भेस पहचानने में इंसानियत
को दो हज़ार साल लगे; तुम्हारे पास बर्फ़ के मैदानों और अंतरिक्ष यानों की
बढ़त है।
\end{remark}

\section{अभ्यास}

\begin{exercise}[$\star$]\label{exo:g9:uniform-varied-motion:1}
\omterm{def:g9:uniform-varied-motion:record}{समयबद्ध स्थिति-अभिलेख} क्या है? उसे बनाने के दो तरीक़े बताओ।
\end{exercise}

\begin{exercise}[$\star$]\label{exo:g9:uniform-varied-motion:2}
\qty{0.1}{s} की टिकों वाली एक पट्टी \qty{4.0}{cm}, $4.0$, $4.0$,
$4.0$ के फ़ासले दिखाती है। गति का निदान करो और उसकी \omterm{def:g7:motion-average-speed:formula}{चाल} \unit{m/s} में
बताओ।
\end{exercise}

\begin{exercise}[$\star$]\label{exo:g9:uniform-varied-motion:3}
एक और पट्टी: फ़ासले $1.0$, $2.5$, $4.5$, \qty{7.0}{cm}. निदान करो ---
और पहले तथा आख़िरी अंतराल की चालें निकालो।
\end{exercise}

\begin{exercise}[$\star$]\label{exo:g9:uniform-varied-motion:4}
स्कूल की एक दौड़ पर औसत और \omterm{def:g9:uniform-varied-motion:instant}{तात्क्षणिक चाल} में फ़र्क़ करो: फाटक के पास
रखा रडार कौन-सी पढ़ता है, और परिवार का ``दरवाज़े से दरवाज़े तक बीस
मिनट'' कौन-सी निकालता है?
\end{exercise}

\begin{exercise}[$\star$]\label{exo:g9:uniform-varied-motion:5}
किन गतियों में औसत और \omterm{def:g9:uniform-varied-motion:instant}{तात्क्षणिक चाल} एक हो जाती हैं? वहीं और सिर्फ़
वहीं क्यों?
\end{exercise}

\begin{exercise}[$\star$]\label{exo:g9:uniform-varied-motion:6}
मुक्त रूप से गिरते पत्थर की \omterm{def:g7:motion-average-speed:formula}{चाल} हर सेकंड कितनी बढ़ती है, और वह बढ़त
किस स्थानीय अंक के बराबर होती है?
\end{exercise}

\begin{exercise}[$\star$]\label{exo:g9:uniform-varied-motion:7}
वह शांत वाक्य बताओ। वह उस पिंड के बारे में क्या कहता है जिस पर कोई \omterm{def:g2:pushes-and-pulls:force}{बल}
लगता ही नहीं --- और उड़ते यात्री विमान की चार-तरफ़ा बराबरी के बारे में?
\end{exercise}

\begin{exercise}[$\star$]\label{exo:g9:uniform-varied-motion:8}
``चलते रहने के लिए गति को \omterm{def:g2:pushes-and-pulls:force}{बल} चाहिए।'' इस पुरानी ग़लती को बहते
अंतरिक्ष यान और घर्षण के भेस से दोषी ठहराओ।
\end{exercise}

\begin{exercise}[$\star\star$]\label{exo:g9:uniform-varied-motion:9}
एक पट्टी (\qty{0.1}{s} की टिकें) \omterm{def:g2:measuring-length:units}{सेंटीमीटर} में बताती है: $6.0$,
$5.0$, $4.0$, $3.0$, $2.0$. निदान करो; हर अंतराल की \omterm{def:g7:motion-average-speed:formula}{चाल} निकालो; और
अगर यही नमूना चलता रहे तो पट्टी के भविष्य की भविष्यवाणी करो।
\end{exercise}

\begin{exercise}[$\star\star$]\label{exo:g9:uniform-varied-motion:10}
हल किए गए ठेले ने \qty{0.1}{s} की हर टिक पर \qty{0.1}{m/s} कमाई।
इस बढ़त को प्रति \emph{सेकंड} बताओ --- और ठेले की कमाई की तुलना मुक्त
पतन की कमाई से करो।
\end{exercise}

\begin{exercise}[$\star\star$]\label{exo:g9:uniform-varied-motion:11}
कोई पैराशूट वाला तेज़ और तेज़ गिरता है --- और फिर, छतरी खुलते ही, टिकी
हुई \qty{5}{m/s} पर। दोनों अवस्थाएँ उस शांत वाक्य से पढ़ो: शुरू में
असंतुलित क्या है, और आख़िर तक कौन-सी बराबरी बन चुकी है?
\end{exercise}

\begin{exercise}[$\star\star\star$]\label{exo:g9:uniform-varied-motion:12}
किसी ढलान से छोड़कर कालीन पर आती खिलौना कार पर प्रयोगशाला का पूरा
फ़ैसला डिज़ाइन करो: कौन-सा अभिलेख लेना है, कौन-से अंक निकालने हैं, दो
अंकों वाला अपेक्षित निदान क्या है (कौन-सा अंक ढलान पर, कौन-सा कालीन
पर), और हर अंक का बल-वाक्य वाला पाठ भी।
\end{exercise}

\section{समस्या: पट्टी-कार्यालय}

\begin{problem}[{सप्ताहांत समस्या --- गति-विश्लेषण कार्यालय में एक
दोपहर; चार पट्टियाँ, चार फ़ैसले; निरीक्षक की
तुकबंदी}]\label{pb:g9:uniform-varied-motion:1}
यह कार्यालय स्कूलों, खेल-क्लबों और अदालतों के लिए समयबद्ध अभिलेखों का
विश्लेषण करता है। हर जगह टिकें \qty{0.1}{s} की हैं; और फ़ासले
\omterm{def:g2:measuring-length:units}{सेंटीमीटर} में आते हैं। रूलर तुम्हारे हाथ में है।

\medskip\noindent\textbf{भाग I --- रोज़मर्रा के फ़ैसले।}
\begin{enumerate}
  \item पट्टी A (गलियारे में चलता कोई): $8.0$, $8.0$, $8.0$, $8.0$,
        $8.0$. फ़ैसला और \omterm{def:g7:motion-average-speed:formula}{चाल} --- \unit{m/s} में और \unit{km/h} में।
  \item पट्टी B (किसी धावक की शुरुआत): $2.0$, $4.0$, $6.0$, $8.0$,
        $10.0$. फ़ैसला, पहली तथा आख़िरी \omterm{def:g7:motion-average-speed:formula}{चाल}, और \omterm{def:g7:motion-average-speed:formula}{चाल} की प्रति-सेकंड
        बढ़त।
  \item पट्टी C (लुढ़कती गेंद, जो रेत से मिल जाती है): $10.0$, $7.0$,
        $4.5$, $2.5$, $1.0$. फ़ैसला --- और रेत के बारे में बल-वाक्य
        का बयान।
  \item पट्टी D (एक पहेली): $5.0$, $5.0$, $5.0$, $7.5$, $10.0$. D की
        दो अंकों वाली कहानी सुनाओ, और उस टिक पर निशान लगाओ जहाँ कुछ
        हुआ।
\end{enumerate}

\medskip\noindent\textbf{भाग II --- कार्यालय की भौतिकी मेज़।}
\begin{enumerate}[resume]
  \item एक मुवक्किल अड़ा है कि उसका सामान पहुँचाने वाला स्कूटर
        ``औसतन \qty{30}{km/h} पर था, हुज़ूर''। जबकि रडार ने चौराहे पर
        \qty{55}{km/h} दर्ज किया। मुवक्किल को समझाओ कि दोनों अंक सच
        कैसे हो सकते हैं, और क़ानून कौन-सा पढ़ता है।
  \item पट्टी B का धावक \omterm{def:g7:motion-average-speed:formula}{चाल} बराबर क़दमों से कमाता है। ऐसी गति पर
        कार्यालय कौन-सा एक शब्द मुहर की तरह लगाता है, और उस बढ़त का
        टिका होना धकेलने वाले \omterm{def:g2:pushes-and-pulls:force}{बल} के बारे में क्या इशारा करता है?
  \item बल-वाक्य कहता है कि पट्टी A का गलियारे वाला \omterm{def:g2:balance-equilibrium:equilibrium}{संतुलित} बलों में
        चल रहा है। टिककर चलने वाले के लिए उस बराबरी का नाम बताओ (क्या
        धकेलता है, क्या विरोध करता है)।
  \item किसी पत्थर के गिरने की पट्टी अंतरालों की चालें $0.98$,
        $1.96$, \qty{2.94}{m/s} दिखाती है। प्रति-सेकंड बढ़त की जाँच
        करो और उस अंक का नाम बताओ जिसे वह दोहरा रही है।
\end{enumerate}

\medskip\noindent\textbf{भाग III --- अदालत का मामला।}
एक स्केटबोर्ड ढलान से लुढ़ककर स्कूल के आँगन के पार क्यारी में जा घुसा;
चौकीदार ``लापरवाह \omterm{def:g7:motion-average-speed:formula}{चाल}'' को दोषी ठहराता है, जबकि स्केट चलाने वाला दावा
करता है कि ``मैं तो पूरे समय धीमा ही होता जा रहा था''। आँगन के कैमरे से
फ़ासले मिलते हैं: ढलान से निकलते हुए $12.0$, फिर $11.5$, $11.0$,
$10.5$, और क्यारी में घुसते हुए $10.0$.
\begin{enumerate}[resume]
  \item स्केट चलाने वाले के दावे पर फ़ैसला --- अंतरालों की चालों समेत।
  \item निकलने की \omterm{def:g7:motion-average-speed:formula}{चाल} और क्यारी में घुसने की \omterm{def:g7:motion-average-speed:formula}{चाल} \unit{km/h} में: किसी
        तेज़ साइकिल सवार की \qty{15}{km/h} के मुक़ाबले क्या इनमें से
        कोई ``लापरवाह'' थी?
  \item चौकीदार पूछता है कि बोर्ड समतल ज़मीन पर अपने आप क्यों नहीं रुक
        गया, ``जैसे चीज़ें क़ुदरतन रुक जाती हैं''। कार्यालय का जवाब,
        \omterm{def:g2:pushes-and-pulls:force}{बल} की ज़बान के एक वाक्य में।
  \item फ़ाइल निरीक्षक की तुकबंदी से बंद करो --- दो पंक्तियाँ: बराबर
        फ़ासलों का मतलब क्या है, और बदलते फ़ासले क्या माँगते हैं (उसी
        शांत उपवाक्य के मुताबिक़, एक कारण)।
\end{enumerate}
\end{problem}
